% Paper B: introduction (one idea, running example first).
\section{Introduction}\label{sec:intro}
Consider a quantity that is carried by a flow and diffuses at the same time, observed with noise. A natural model is
\begin{equation}\label{eq:intro-tpd}
  (\partial_t+c\,\partial_x)(\partial_t-\nu\,\partial_x^2)\,u=0,
\end{equation}
whose solutions are sums of a travelling profile and a diffusing one. A method that learns equations by regressing
$u_t$ on a library of terms such as $u_x$, $u_{xx}$, $u_{xxx}$ looks for one equation of the form $u_t=\sum_k\beta_k\,\partial_x^ku$; the
data in fact obey the product of two equations, and no single sum of terms fits them. On our benchmark, every
library-regression method we tested returns a wrong equation for \eqref{eq:intro-tpd} at $5\%$ noise, without warning.

This paper argues that equation learning should group the data by the \emph{factors} of the equation, and develops that
idea from the algebra to an algorithm that quantifies its uncertainty. In Fourier variables a constant-coefficient operator
is a polynomial in $(\tau,\xi)$; the operator in \eqref{eq:intro-tpd} is $(\tau+c\,\xi)(\tau-\nu\,\xi^2)$. Every Fourier mode
$e^{\lambda t+i\kappa x}$ of a solution lies on the zero set of one factor: the modes of the first factor on the line
$\lambda=-ic\kappa$, those of the second on the parabola $\lambda=-\nu\kappa^2$ (\cref{fig:branches}). Learning the equation
means learning these groups.

% Shared figure for the introductions: the two groups of the running example as curves of modes.
\begin{figure}[tbp]
\centering
\begin{tikzpicture}[x=0.40cm,y=0.40cm,>=Stealth,font=\footnotesize]
  % (a) oscillation slice
  \begin{scope}
    \draw[->,gray] (-4.5,0)--(4.8,0) node[right,black]{$\kappa$};
    \draw[->,gray] (0,-4.5)--(0,4.8) node[above,black]{$\operatorname{Im}\lambda$};
    \draw[very thick,color={rgb,255:red,42;green,120;blue,214}] (-4.2,4.2)--(4.2,-4.2);
    \draw[very thick,color={rgb,255:red,235;green,104;blue,52},dashed] (-4.2,0)--(4.2,0);
    \foreach \k in {-3,-2,-1,1,2,3} \fill[color={rgb,255:red,42;green,120;blue,214}] (\k,-\k) circle (2.6pt);
    \foreach \k in {-2,-1,1,2} \fill[color={rgb,255:red,235;green,104;blue,52}] (\k,0) +(-2.4pt,-2.4pt) rectangle +(2.4pt,2.4pt);
    \node[anchor=west,align=left] at (1.2,3.3) {transport\\$\lambda=-ic\kappa$};
    \node at (0,-5.8) {(a) oscillation};
  \end{scope}
  % (b) damping slice
  \begin{scope}[xshift=10.6cm]
    \draw[->,gray] (-4.5,0)--(4.8,0) node[right,black]{$\kappa$};
    \draw[->,gray] (0,-4.5)--(0,1.6) node[above,black]{$\operatorname{Re}\lambda$};
    \draw[very thick,color={rgb,255:red,42;green,120;blue,214},dashed] (-4.2,0)--(4.2,0);
    \draw[very thick,color={rgb,255:red,235;green,104;blue,52},domain=-3.8:3.8,samples=60] plot (\x,{-0.3*\x*\x});
    \foreach \k in {-3,-2,-1,1,2,3} \fill[color={rgb,255:red,42;green,120;blue,214}] (\k,0) circle (2.6pt);
    \foreach \k in {-2,-1,1,2} \fill[color={rgb,255:red,235;green,104;blue,52}] (\k,{-0.3*\k*\k}) +(-2.4pt,-2.4pt) rectangle +(2.4pt,2.4pt);
    \node[anchor=west,align=left] at (1.8,1.6) {diffusion\\$\lambda=-\nu\kappa^2$};
    \node at (0,-5.8) {(b) damping};
  \end{scope}
\end{tikzpicture}
\caption{The idea of the paper on the equation $(\partial_t+c\,\partial_x)(\partial_t-\nu\,\partial_x^2)u=0$. Every Fourier mode
$e^{\lambda t+i\kappa x}$ of a solution lies on the zero set of one factor of the symbol $(\tau+c\xi)(\tau-\nu\xi^2)$: circles
on the line of the transport factor, squares on the parabola of the diffusion factor (schematic; \cref{fig:directions}
shows the modes found in noisy data). Learning the equation means learning these two groups.}\label{fig:branches}
\end{figure}


\paragraph{The grouping principle} We prove that the groups are exactly the factors of the symbol: every solution splits
into one component per factor, uniquely up to an overlap whose dimension is the number of intersection points of their
zero sets (\cref{thm:factor}). Additive terms do not form groups: for two terms whose varieties meet in finitely many points, no solution outside a
finite-dimensional set is a sum of single-term pieces (\cref{prop:terms}), and the interaction between terms cannot be
removed stably (\cref{prop:discont}). The split is stable exactly when the conditioning constant $\gN$ of the sum operator
is positive; for two factors this means a positive Friedrichs angle between their solution spaces, and the error
amplification is $\gN^{-1/2}$ (\cref{thm:hm}). Finally, without a priori bounds, finite data can support a factored
equation but never certify it (\cref{thm:ident,prop:certify}).

\paragraph{Learning the groups} The algorithm groups twice (\cref{fig:roadmap}). Within one posterior sample of the field
it reads the temporal eigenvalues of every wavenumber with a matrix pencil and assigns the resulting modes to factors by
fitting the zero sets directly, a clustering problem whose clusters are curves. Across posterior samples it clusters the
factors themselves with a Dirichlet-process mixture. Every decision is then checked against the raw observations, and
the output is a distribution over equations: factors with $90\%$ intervals, the probability of each number of groups,
the energy share and separation of each group, and a label that says whether the data support the equation, support
it only within the chosen class, reject it, or cannot decide.

\paragraph{Contributions}
\begin{enumerate}
  \item The grouping principle and its limits: \cref{thm:factor} with the exact ambiguity $\deltaP$, the negative results
  for terms (\cref{prop:terms,prop:discont}), the conditioning of the split (\cref{thm:hm}), its Bayesian counterpart
  (\cref{prop:bayes-angle}), and a counting condition on how many wavenumbers a group needs (\cref{prop:counting}).
  \item An algorithm that learns the groups from noisy, scattered data, with a sampling-based treatment of the
  high-frequency band and labels that encode what finite data can support.
  \item Evidence on twenty random equations not used in development ($19$ right or equivalent, $14$ of them with a confident
  label; no confident wrong answer, against $13$ or more for each of the three library baselines run on this set), on
  $60$ grid inputs with independent fields and noise (all right; weak SINDy is wrong without warning on $13$) and on ten benchmark equations in three designs, on which the
  parameters were tuned ($30$ of $30$ factor sets recovered; the best library regression is wrong without warning on $5$).
\end{enumerate}

% Shared: relation to prior work (what is classical, what is new).
\paragraph{Relation to prior work}
Most methods that learn a PDE from data regress a time derivative on a library of candidate terms: SINDy and
PDE-FIND~\cite{BruProKut16,RudBruProKut17}, their weak, ensemble and Bayesian variants~\cite{MesBor21,FasKutBru22,HirBarKut22},
and Gaussian-process formulations of linear operators~\cite{RaiPerKar17}. Their unit is the \emph{term}. Dynamic mode
decomposition~\cite{Sch10,SasKut22} works with eigenvalues, as we do, but returns a linear operator on the data, not an
equation, and does not group eigenvalues by the factor they belong to. Exponential fitting by matrix
pencils~\cite{HuaSar90} is a classical tool of signal processing; the algebra of constant-coefficient operators, the
Ehrenpreis--Palamodov fundamental principle and the solvability results behind \cref{thm:factor} are classical
analysis~\cite{Ehr70,Pal70,Mal56,Hor90,Obe90}; Dirichlet-process mixtures are standard Bayesian
clustering~\cite{BleJor06}. What is new here is the connection between them. (i) We identify the \emph{factor} of the
symbol, not the term, as the unit a learning method should group by, and we prove that, for two terms whose varieties meet in finitely many points, per-term grouping is impossible
outside a finite-dimensional set of solutions (\cref{prop:terms,prop:discont}). (ii) We state what finite data can and
cannot certify about such groups (\cref{thm:ident,prop:certify}) and give the exact conditioning of the
separation between groups (\cref{thm:hm,prop:bayes-angle}). (iii) We build an algorithm that groups twice, modes into
factors within a sample and factors across posterior samples, and whose output carries these limits as labels and
intervals.


\paragraph{Scope and organization} Everything proved and implemented concerns scalar linear equations with constant
coefficients in one time and one space variable, observed on a rectangle. \Cref{sec:theory} proves the grouping principle
and its limits. \Cref{sec:within} describes how the groups are
learned within one sample and \cref{sec:across} how they are grouped across samples. \Cref{sec:evidence} follows one data
set through every stage and reports benchmarks, and \cref{sec:discussion} discusses limits, what is lost beyond the linear ring, and extensions. \Cref{eq:intro-tpd}
with $c=1$, $\nu=0.3$ is the running example throughout; \cref{tab:notation} collects the notation and proofs are in
\cref{app:proofs}. \emph{Reader's guide.} The algorithm uses four results: the factor decomposition (\cref{thm:factor}), the
stability of the split (\cref{thm:hm}), the counting condition (\cref{prop:counting}) and the limits of finite data
(\cref{prop:certify}), which explain the labels. A reader interested mainly in the method can skim the material on
exponential-polynomials and identification on a first reading.

% Shared figure: roadmap of the method, each box with the section and the result it rests on.
\begin{figure}[tbp]
\centering
\begin{tikzpicture}[
  node distance=3.5mm and 4.5mm,
  box/.style={draw=black!55, rounded corners=2pt, align=center, font=\footnotesize, inner sep=3pt,
              minimum height=10mm, text width=20mm, fill=white},
  grp/.style={box, fill=blue!6},
  thy/.style={font=\scriptsize\itshape, text=black!60, align=center},
  arr/.style={-{Stealth[length=2mm]}, black!60, thick}]
\node[box] (data) {noisy data\\$\{(t_i,x_i,y_i)\}$};
\node[box, right=of data] (gp) {field samples\\$u^{(1)},\dots,u^{(\Nsamp)}$};
\node[box, right=of gp] (modes) {modes $(\kappa,\lambda)$\\per wavenumber};
\node[grp, right=of modes] (fac) {factors $Q_j$\\(groups of modes)};
\node[grp, below=8mm of fac] (cloud) {cloud of factors\\over samples};
\node[box, left=of cloud] (clus) {clusters and\\$P(k)$, $P(\Mfit)$};
\node[box, left=of clus] (split) {components $u_j$,\\$w_j$, $\gN$, $\rho$};
\node[box, left=of split] (out) {equation with\\intervals and label};
\draw[arr] (data)--(gp); \draw[arr] (gp)--(modes); \draw[arr] (modes)--(fac);
\draw[arr] (fac)--(cloud); \draw[arr] (cloud)--(clus); \draw[arr] (clus)--(split); \draw[arr] (split)--(out);
\node[thy, above=1mm of gp] {\S\ref{sec:alg-interp}--\ref{sec:nyquist}};
\node[thy, above=1mm of modes] {\S\ref{sec:alg-pencil}, Lem.~\ref{lem:ep}};
\node[thy, above=1mm of fac] {\S\ref{sec:alg-structfit}, Thm.~\ref{thm:factor}};
\node[thy, below=1mm of cloud] {\S\ref{sec:bayes-cloud}};
\node[thy, below=1mm of clus] {\S\ref{sec:alg-loop}, \S\ref{sec:bayes-cloud}};
\node[thy, below=1mm of split] {Thms.~\ref{thm:hm}, \ref{thm:eint}};
\node[thy, below=1mm of out] {\S\ref{sec:alg-labels}, Thm.~\ref{thm:ident}};
\end{tikzpicture}
\caption{The method at a glance. Shaded: the two grouping steps, of modes into factors within one sample and of
factors across samples. Each box names the section that describes it and the result that justifies it.}\label{fig:roadmap}
\end{figure}

% Shared module: notation table (one symbol, one meaning). Checked by tools/check_notation.py.
\begin{table}[htbp]
\centering\small
\caption{Notation for the analysis. Each symbol has one meaning throughout the paper; \cref{tab:notation-alg} lists the notation of the algorithm.}\label{tab:notation}
\begin{tabularx}{\linewidth}{@{}l X l@{}}
\toprule
Symbol & Meaning & Introduced\\
\midrule
\multicolumn{3}{@{}l}{\emph{Operators and their algebra}}\\
$\ring=\C[\tau,\xi]$ & ring of constant-coefficient operators; $\tau\leftrightarrow\partial_t$, $\xi\leftrightarrow\partial_x$ & \cref{def:ring}\\
$P,\ Q_j,\ P_i,\ e_i$ & an operator and its symbol; a factor; an irreducible factor and its multiplicity & \cref{sec:ring,sec:factor}\\
$\Var(P)$ & characteristic variety $\{P=0\}\subset\C^2$ & \cref{def:ring}\\
$\kappa,\ \lambda(\kappa)$ & real wavenumber ($\xi=i\kappa$); temporal rate of a mode & \cref{sec:ring}\\
$\Lambda_j(\xi)$ & branch of an evolution factor $Q_j=\tau-\Lambda_j(\xi)$ & \eqref{eq:lagrange}\\
$\Ann(u)$ & annihilator: all operators that kill $u$ & \cref{def:ann}\\
$\Gu,\ \Ju,\ \sigp(u)$ & intrinsic operator $\gcdop\Ann(u)$; cofactor ideal; point spectrum $\Var(\Ju)$ & \eqref{eq:GJ}\\
$\EP$ & exponential-polynomials & \cref{def:ep}\\
$N,\ \deltaP$ & number of factors; dimension of the overlap between their solution spaces & \cref{thm:factor}\\
$m_j,\ d_j$ & degree of factor $j$ in $\tau$ and in $\xi$ & \cref{sec:alg-loop}\\
\multicolumn{3}{@{}l}{\emph{Geometry and energy}}\\
$\sumop$ & sum operator $(u_1,\dots,u_N)\mapsto\sum_ju_j$ & \cref{sec:cond}\\
$\gN$ & conditioning constant of $\sumop$; error amplification $\gN^{-1/2}$ & \eqref{eq:gammaN}\\
$\cF,\ \thetaF$ & Friedrichs cosine and angle; $\gamma_2=1-\cF$ & \cref{sec:cond}\\
$\Hs$ & Sobolev space of index $s$ & \cref{sec:cond}\\
$K;\ \mathcal A_n,\ \flow_n,\ \Delta_n$ & number of additive terms of a factor; a term; its flow $e^{t\mathcal A_n}$; increment $\flow_n-\Id$ & \cref{sec:eint}\\
$\hat{\mathcal A}_n(\kappa)$ & Fourier symbol of the term $\mathcal A_n$ & \eqref{eq:eint-fourier}\\
$\Eint,\ \rho$ & interaction term; interaction ratio $\norm{\Eint}/\norm{u(T)-u_0}$ & \eqref{eq:eint-def}, \eqref{eq:energy}\\
$w_j$ & energy share of group $j$ & \eqref{eq:energy}\\
\bottomrule
\end{tabularx}
\end{table}

\begin{table}[htbp]
\centering\small
\caption{Notation for the algorithm (continued from \cref{tab:notation}).}\label{tab:notation-alg}
\begin{tabularx}{\linewidth}{@{}l X l@{}}
\toprule
Symbol & Meaning & Introduced\\
\midrule
$T,\ L$ & length of the time window; spatial period & \cref{sec:alg-interp}\\
$\Delta t,\ \Delta x;\ h$ & grid spacings; mean spacing $\sqrt{TL/n}$ of $n$ scattered points & \cref{sec:nyquist}\\
$\noise,\ \sigbin$ & observation noise standard deviation; its value in one Fourier bin & \cref{sec:alg-interp,sec:nyquist}\\
$\mathcal K,\ z$ & band of usable wavenumbers; detection threshold in units of $\sigbin$ & \cref{sec:nyquist}\\
$h_\kappa$ & spread $\max_{\ell,\ell'}\abs{\kappa_\ell-\kappa_{\ell'}}$ of the wavenumbers on which a factor is fitted & \cref{prop:counting}\\
$A_p,\ \sigma_{\lambda,p}$ & amplitude of mode $p$; standard deviation of its rate & \cref{sec:alg-pencil}\\
$\Nsamp,\ u^{(q)},\ \pi_q$ & number of posterior field samples; sample $q$; its weight & \cref{sec:alg-interp}\\
$\Mfit$ & a structure (number and degrees of the factors) & \cref{sec:alg-loop}\\
$B_j$ & orthonormal basis of the sampled kernel of $Q_j$ & \cref{sec:alg-interference}\\
$E_j$ & kernel exponentials of $Q_j$ evaluated at the observation points & \cref{sec:alg-accept}\\
$\omega^2$ & prior variance of the component coefficients & \eqref{eq:linmodel}\\
$\eta_r,\ \ell^{(r)}$ & amplitudes and length scales of the Gaussian-process kernel & \eqref{eq:gp}\\
$\eta,\ \eta_h,\ \pi_h,\ N_\eta$ & hyperparameters; a draw of them; its weight; number of draws & \eqref{eq:marginal-eta}\\
$\beta,\ \cloud$ & coefficient vector of a factor; cloud of factors & \eqref{eq:cloud}\\
$\alpha$ & concentration of the Dirichlet process & \cref{sec:bayes-cloud}\\
$k,\ \keff,\ \bar w_j$ & number of clusters (groups) found; effective number of groups; cluster masses & \cref{sec:bayes-cloud}\\
\FC, \IDK, \MIS, \UND & output labels & \cref{sec:alg-labels}\\
\bottomrule
\end{tabularx}
\end{table}

